\chapter{کد}

بلوک‌های کد با \lr{minted} و پایتون‌های \lr{Pygments} رنگ‌آمیزی می‌شوند و در
کادری با نوار بالایی می‌نشینند که نام زبان و در صورت وجود نام فایل را نشان
می‌دهد. کد همیشه از چپ به راست می‌ماند، حتی وسط متن فارسی. برای کامپایل به
\lr{\texttt{-shell-escape}} نیاز دارید.

\section{بلوکی با نام فایل}

\begin{code}{python}{fib.py}
from functools import lru_cache

@lru_cache(maxsize=None)
def fib(n: int) -> int:
    """Nothing clever. The cache does the work."""
    return n if n < 2 else fib(n - 1) + fib(n - 2)

print([fib(n) for n in range(10)])
\end{code}

اگر قطعه کد به فایلی تعلق ندارد، آرگومان دوم را خالی بگذارید:

\begin{code}{rust}{}
fn main() {
    let total: u32 = (1..=100).filter(|n| n % 3 == 0 || n % 5 == 0).sum();
    println!("{total}");
}
\end{code}

\section{زبان‌های دیگر}

\begin{code}{sql}{report.sql}
select
  date_trunc('month', created_at) as month,
  count(*) filter (where status = 'shipped') as shipped
from orders
where created_at >= now() - interval '1 year'
group by 1
order by 1;
\end{code}

\section{ترمینال}

خروجی پوسته روی زمینه‌ی تیره می‌نشیند، چون ترمینال همین شکلی است و وانمود
کردن به چیز دیگر به کسی کمک نمی‌کند.

\begin{terminal}
$ latexmk -xelatex -shell-escape main.tex
Latexmk: Run number 1 of rule 'xelatex'
Output written on main.pdf (24 pages, 412918 bytes).
Latexmk: All targets (main.pdf) are up-to-date
\end{terminal}

\section{متن بدون رنگ‌آمیزی}

برای خروجی، لاگ، یا هرچه که \lr{Pygments} فقط خرابش می‌کند:

\begin{plaincode}
notebook-template/
├── styles/     فایل‌های سبک، هر کدام یک وظیفه
├── fonts/      XM Yekan, Inter, SF Mono
├── vars/       رنگ‌ها و مشخصات — این‌ها را ویرایش کنید
├── sections/   نوشته‌ی خودتان
└── main.tex
\end{plaincode}

\section{درون‌خطی}

برای پاک‌کردن فایل‌های کمکی \cmd{latexmk -c} و برای حذف خود \lr{PDF}
\cmd{latexmk -C} را اجرا کنید. حافظه‌ی نهان \lr{minted} در
\cmd{\_minted-main/} می‌ماند و از پیش در \lr{\texttt{.gitignore}} هست.

\begin{tip}[گنجاندن فایل واقعی]
\lr{\texttt{\textbackslash inputcode\{python\}\{src/model.py\}}} فایل را
مستقیم از دیسک می‌خواند، پس نوشته هیچ‌وقت از کدی که توضیحش می‌دهد عقب
نمی‌ماند.
\end{tip}

\begin{note}[تغییر شِمای رنگ]
شِمای رنگ‌آمیزی یک‌جا در \lr{\texttt{styles/code.sty}} تنظیم می‌شود
(\lr{\texttt{style = friendly}}). با \cmd{pygmentize -L styles} بقیه را
ببینید؛ \lr{\texttt{tango}} و \lr{\texttt{xcode}} روی کاغذ خوب چاپ می‌شوند.
\end{note}
